
%(BEGIN_QUESTION)
% Copyright 2006, Tony R. Kuphaldt, released under the Creative Commons Attribution License (v 1.0)
% This means you may do almost anything with this work of mine, so long as you give me proper credit

Anta at vi måler strømningsraten til en svak syreløsning med en magnetisk flowmåler. Syrenes ledningsevne ligger godt innenfor akseptabelt område for denne måleren, og den fungerer derfor helt fint.

Anta nå at syreløsningen blir sterkere (høyere syrekonsentrasjon). Dette vil øke ledningsevnen i løsningen, fordi det nå finnes flere ioner som kan bære elektrisk strøm. Hvilken effekt vil dette ha på kalibreringen av den magnetiske flowmåleren? Må noen rekalibrere flowmåleren for at den igjen skal måle syrestrømmen korrekt? Hvis ja, blir dette et nullpunkts- eller spansskift? I hvilken retning vil nullpunkt og/eller span forskyves, opp eller ned? Forklar svaret/svarene dine.

\underbar{file i00729}
%(END_QUESTION)





%(BEGIN_ANSWER)

Endringer i væskens ledningsevne har neglisjerbar effekt på flowmålerens kalibrering.

\vskip 10pt

En vanlig misforståelse om magnetiske flowmålere gjelder sammenhengen mellom væskens ledningsevne og flowmålerens kalibrering. Så lenge ledningsevnen holder seg innenfor målerens akseptable område, har endringer i ledningsevnen neglisjerbar effekt på kalibreringen. Målerens spenningsmålende krets har så mye høyere impedans enn den elektriske banen gjennom væsken, at endringer i væskens ledningsevne blir «druknet» av målerens langt høyere inngangsimpedans.

%(END_ANSWER)





%(BEGIN_NOTES)


%INDEX% Måling, flow: magnetisk

%(END_NOTES)


